\SourceLeaf{77}%
\begin{center}
{\large THEORIA RESIDUORUM BIQUADRATICORUM.\par}
\bigskip COMMENTATIO PRIMA.
\end{center}
\BookRule{28mm}

\Article{1}
Theoria residuorum quadraticorum ad pauca theoremata fundamentalia reducitur, pulcherrimis Arithmeticae Sublimioris cimeliis adnumeranda, quae primo per inductionem facile detecta, ac dein multifariis modis ita demonstrata esse constat, ut nihil amplius desiderandum relictum sit.

Longe vero altioris indaginis est theoria residuorum cubicorum et biquadraticorum. Quam quum inde ab anno 1805 perscrutari coepissemus, praeter ea, quae quasi in limine sunt posita, nonnulla quidem theoremata specialia se obtulerunt, tum propter simplicitatem suam, tum propter demonstrationum difficultatem valde insignia: mox vero comperimus, principia Arithmeticae hactenus usitata ad theoriam generalem stabiliendam neutiquam sufficere, quin potius hanc necessario postulare, ut campus Arithmeticae Sublimioris infinities quasi promoveatur, quod quomodo intelligendum sit, in continuatione harum disquisitionum clarissime elucebit. Quamprimum hunc campum novum ingressi sumus, aditus ad cognitionem theorematum simplicissimorum totam theoriam exhaurientium per inductionem statim patuit: sed ipsorum demonstrationes tam profunde latuerunt, ut post multa demum tentamina irrita tandem in lucem protrahi potuerint.

Quum iam ad promulgationem harum lucubrationum accingamur, a theoria residuorum biquadraticorum initium faciemus, et quidem in hac prima commen%
\SourceLeaf{78}%
tatione disquisitiones eas explicabimus, quas iam cis campum Arithmeticae ampliatum absolvere licuit, quae illuc viam quasi sternunt, simulque theoriae divisionis circuli quaedam nova incrementa adiungunt.

\Article{2}
Notionem residui biquadratici in \emph{Disquisitionibus Arithmeticis} art. 115 introduximus: scilicet numerus integer $a$, positivus seu negativus, integri $p$ residuum biquadraticum vocatur, si $a$ secundum modulum $p$ biquadrato congruus fieri potest, et perinde non-residuum biquadraticum, si talis congruentia non exstat. In omnibus disquisitionibus sequentibus, ubi contrarium expressis verbis non monetur, modulum $p$ esse numerum primum (imparem positivum) supponemus, atque $a$ per $p$ non divisibilem, quum omnes casus reliqui ad hunc facillime reduci possint.

\Article{3}
Manifestum est, omne residuum biquadraticum numeri $p$ eiusdem quoque residuum quadraticum esse, et proin omne non-residuum quadraticum etiam non-residuum biquadraticum. Hanc propositionem etiam convertere licet, quoties $p$ est numerus primus formae $4n+3$. Nam si in hoc casu $a$ est residuum quadraticum ipsius $p$, statuamus $a\equiv bb\; (\mathrm{mod.}\,p)$, ubi $b$ vel residuum quadraticum ipsius $p$ erit vel non-residuum: in casu priori statuemus $b\equiv cc$, unde $a\equiv c^4$, i. e. $a$ erit residuum biquadraticum ipsius $p$; in casu posteriori $-b$ fiet residuum quadraticum ipsius $p$ (quoniam $-1$ est non-residuum cuiusvis numeri primi formae $4n+3$), faciendoque $-b\equiv cc$, erit ut antea $a\equiv c^4$, atque $a$ residuum biquadraticum ipsius $p$. Simul facile perspicietur, alias solutiones congruentiae $x^4\equiv a\; (\mathrm{mod.}\,p)$, praeter has duas $x\equiv c$ et $x\equiv-c$ in hoc casu non dari. Quum hae propositiones obviae integram residuorum biquadraticorum theoriam pro modulis primis formae $4n+3$ exhauriant, tales modulos a disquisitione nostra omnino excludemus, sive hanc ad modulos primos formae $4n+1$ limitabimus.

\Article{4}
Existente itaque $p$ numero primo formae $4n+1$, propositionem art. praec. convertere non licet: nempe exstare possunt residua quadratica, quae non sunt simul residua biquadratica, quod evenit, quoties residuum quadraticum congruum est quadrato non-residui quadratici. Statuendo enim $a\equiv bb$, existente $b$ non-%
\SourceLeaf{79}%
residuo quadratico ipsius $p$, si congruentiae $x^4\equiv a$ satisfieri posset, per valorem $x\equiv c$, foret $c^4\equiv bb$, sive productum $(cc-b)(cc+b)$ per $p$ divisibile, unde $p$ vel factorem $cc-b$ vel alterum $cc+b$ metiri deberet, i. e. vel $+b$ vel $-b$ foret residuum quadraticum ipsius $p$, et proin uterque (quoniam $-1$ est residuum quadraticum), contra hyp.

Omnes itaque numeri integri per $p$ non divisibiles in tres classes distribui possent, quarum prima contineat residua biquadratica, secunda non-residua biquadratica ea, quae simul sunt residua quadratica, tertia non-residua quadratica. Manifesto sufficit, tali classificationi solos numeros $1,2,3\dotsfive p-1$ subiicere, quorum semissis ad classem tertiam reduceretur, dum altera semissis inter classem primam et secundam distribueretur.

\Article{5}
Sed praestabit, quatuor classes stabilire, quarum indoles ita se habeat.

Sit $A$ complexus omnium residuorum biquadraticorum ipsius $p$, inter $1$ et $p-1$ (inclus.) sitorum, atque $e$ non-residuum quadraticum ipsius $p$ ad arbitrium electum. Sit porro $B$ complexus residuorum minimorum positivorum e productis $eA$ secundum modulum $p$ oriundorum, et perinde $C,D$ resp. complexus residuorum minimorum positivorum e productis $eeA$, $e^3A$ secundum modulum $p$ prodeuntium. His ita factis facile perspicitur, singulos numeros $B$ inter se diversos fore, et perinde singulos $C$, nec non singulos $D$; cifram autem inter omnes hos numeros occurrere non posse. Porro patet, omnes numeros, in $A$ et $C$ contentos, esse residua quadratica ipsius $p$, omnes autem in $B$ et $D$ non-residua quadratica, ita ut certe complexus $A,C$ nullum numerum cum complexu $B$ vel $D$ communem habere possint. Sed etiam neque $A$ cum $C$, neque $B$ cum $D$ ullum numerum communem habere potest. Supponamus enim

I. numerum aliquem ex $A$, e. g. $a$ etiam in $C$ inveniri, ubi prodierit e producto $eea'$ ipsi congruo, existente $a'$ numero e complexu $A$. Statuatur $a\equiv\alpha^4$, $a'\equiv\alpha'^4$; accipiaturque integer $\Theta$ ita, ut fiat $\Theta\alpha'\equiv1$. His ita factis erit $ee\alpha'^4\equiv\alpha^4$, adeoque multiplicando per $\Theta^4$,

\[
ee\equiv\alpha^4\Theta^4
\]

i. e. $ee$ residuum biquadraticum, adeoque $e$ residuum quadraticum, contra hyp.

\SourceLeaf{80}%
II. Perinde supponendo, aliquem numerum complexibus $B,D$ communem esse, atque e productis $ea$, $e^3a'$ prodiisse, existentibus $a,a'$ numeris e complexu $A$, e congruentia $ea\equiv e^3a'$ sequeretur $a\equiv eea'$, adeoque haberetur numerus, qui e producto $eea'$ oriundus ad $C$ simulque ad $A$ pertineret, quod impossibile esse modo demonstravimus.

Porro facile demonstratur, \emph{omnia} residua quadratica ipsius $p$, inter $1$ et $p-1$ incl. sita, necessario vel in $A$ vel in $C$, omniaque non-residua quadratica ipsius $p$ inter illos limites necessario vel in $B$ vel in $D$ occurrere debere. Nam

I. Omne tale residuum quadraticum, quod simul est residuum biqua\-draticum, per hyp. in $A$ invenitur.

II. Residuum quadraticum $h$ (ipso $p$ minus), quod simul est non-residuum biquadraticum, statuatur $\equiv gg$, ubi $g$ erit non-residuum quadraticum. Accipiatur integer $\gamma$ talis, ut fiat $e\gamma\equiv g$, eritque $\gamma$ residuum quadraticum ipsius $p$, quod statuemus $\equiv kk$. Hinc erit

\[
h\equiv gg\equiv ee\gamma\gamma\equiv eek^4
\]

Quare quum residuum minimum ipsius $k^4$ inveniatur in $A$, numerus $h$, quippe qui ex illius producto per $ee$ oritur, necessario in $C$ contentus erit.

III. Designante $h$ non-residuum quadraticum ipsius $p$ inter limites $1$ et $p-1$, eruatur inter eosdem limites numerus integer $g$ talis, ut habeatur $eg\equiv h$. Erit itaque $g$ residuum quadraticum, et proin vel in $A$ vel in $C$ contentus: in casu priori $h$ manifesto inter numeros $B$, in posteriori autem inter numeros $D$ invenietur.

Ex his omnibus colligitur, cunctos numeros $1,2,3\dotsfive p-1$ inter quatuor series $A,B,C,D$ ita distribui, ut quivis illorum in una harum reperiatur, unde singulae series $\tfrac14(p-1)$ numeros continere debent. In hac classificatione classes $A$ et $C$ quidem numeros suos essentialiter possident, sed distinctio inter classes $B$ et $D$ eatenus arbitraria est, quatenus ab electione numeri $e$ pendet, qui ipse semper ad $B$ referendus est; quapropter si eius loco alius e classe $D$ adoptatur, classes $B,D$ inter se permutantur.

\Article{6}
Quum $-1$ sit residuum quadraticum ipsius $p$, statuamus, $-1\equiv ff\; (\mathrm{mod.}\,p)$, unde quatuor radices congruentiae $x^4\equiv1$ erunt $1,f,-1,-f$. Quodsi itaque\space%
